\documentclass[11pt,a4paper]{ctexart}
\usepackage[margin=2.1cm]{geometry}
\usepackage{amsmath,amssymb,booktabs,float,graphicx,xcolor,hyperref,enumitem,longtable}
\hypersetup{colorlinks=true,linkcolor=blue!50!black,urlcolor=blue!50!black}
\setlist{itemsep=2pt,topsep=3pt}
\title{H50：自动代数多场约化（AMFR）——锁定实验}
\author{温州肯恩大学 \quad Nonlinear Poisson 项目}
\date{2026 年 11 月 5 日}

\begin{document}
\maketitle

\begin{abstract}
报告 12 的结论是：在 SC26\_01 这类多场矩阵上，能赢 BoomerAMG 的不是减少迭代，而是利用块结构让每次迭代更便宜；建议把 bt-v、bt-j 和整体 AMG 统一成一个只看矩阵的自动方法。
本报告实现了这个方法（AMFR）：输入只有矩阵和场的个数，按三条代数判据自动决定哪些场用对角消去、哪些场联合求解、按什么顺序求解。
在预先锁定的 58 个新算例（自建三温离散 22 个，多群辐射扩散 + 物质温度 36 个）和真实 SC26\_01 上测试：
\textbf{(1)} 在三温问题上，AMFR 自动重现了 bt-v / bt-j，是否消去电子的判断 20/20 与预测一致，总时间是 BoomerAMG 的 0.78--1.19 倍（中位数 0.93）；真实 SC26\_01 上自动选出 bt-v，16.2\,s 对 18.1\,s。
\textbf{(2)} 在多群问题上，BoomerAMG 弱的算例 AMFR 快很多（467 步 → 40 步，0.11 倍；BoomerAMG 500 步不收敛 → 8 步，0.02 倍），收益几乎全部来自“对角消去”。
\textbf{(3)} 但有一类系统性失败：物质温度有热传导、热容小时，AMFR 把它和辐射群分开解，迭代是 BoomerAMG 的 1.7--4.9 倍（时间 1.67--2.25 倍）。
6 条预测中 2 条成立（真实矩阵、结构判断），4 条不成立（稳健性差一个算例、时间中位数、规则质量、对 hypre systems AMG）。
锁定后的诊断（探索性，未锁定）找到了失败原因：耦合判据里对单边耦合做了截断 $\min(1,\cdot)$，物质温度行不是对角占优（耦合/行和 $\approx 2.5$），截断把强耦合判成了弱耦合；这个截断也破坏了判据本应有的单位不变性。
去掉截断并改用每个单元耦合矩阵的谱半径后，这类算例的迭代数降到与 BoomerAMG 相同（155 → 49 步），但时间只是持平，没有变快。
\end{abstract}

\section{方法（锁定前固定）}
输入：稀疏矩阵 $A$，以及 $k$ 个等长、按单元对齐的场（通常的分块多物理排列）。不使用任何问题相关信息。
\begin{enumerate}
\item \textbf{消去。} 若场 $e$ 不是最后一个场、与其他场只有点耦合（对角块）、且非对角占比
\[s_e=q_{95}\Big(\sum_{j\ne c}|a^{ee}_{cj}|\,/\,|a^{ee}_{cc}|\Big)\le\tau_d=2\times10^{-2},\]
就用它的对角消去（$s_e$ 是用对角代替 $A_{ee}$ 的收缩因子上界）。
每次消一个（占比最小的先消），形成 Schur 补 $S_{ij}=A_{ij}-\mathrm{diag}(a_{ie}a_{ej}/d_e)$，对角用无抵消形式
$S_{ii}=(a_{ii}+a_{ie})-a_{ie}(d_e+a_{ei})/d_e$，然后重新计算。
\item \textbf{合并。} 对剩下的场，常数模态（最慢模态）的单元回路增益
$g_{ij}=\min(1,b_{ij}/\rho_i)\min(1,b_{ji}/\rho_j)$，$b_{ij}$ 是 $|S_{ij}|$ 的行和，$\rho_i$ 是 $S_{ii}$ 行和的绝对值。
若 $g_{ij}>\tau_g=0.7$ 的单元比例 $\ge\varphi=10^{-3}$，就把 $i,j$ 合并、联合求解。
\item \textbf{排序。} 块下三角顺序：每次先解受其余未解组影响最小的组，丢掉对后面组的耦合。
\item \textbf{求解。} 每组用一次 BoomerAMG V 循环（与 amg-3d 相同参数），组内按场分块当作一个标量矩阵。
\end{enumerate}
bt-v（消去电子，离子 → 辐射）、bt-j（电子+离子联合 → 辐射）和整体 AMG（全部合并）都是这个规则的特例。单元测试：在 gen3t 基准算例上与 BTElim 的输出相差 $5\times10^{-16}$，在 $G=1$ 上与 BTJoint 完全相同；不丢耦合时为精确逆（$3\times10^{-16}$）。

\paragraph{开发数据与锁定。} 阈值只在已见过的数据上设定：H43/H44 的 20 个 gen3t 区间，以及 8 个多群开发算例（种子 1，$m=32$，3 群）。
$\tau_g$ 原为 0.9，开发数据上 gen3t $F=0.1,G=1$（电子--离子增益 0.89，未合并，31 步对 17 步）失败，改为 0.7；试过对联合块按场的对角中位数缩放，迭代变为 3 倍（36 对 15），已删除。
方案 \texttt{protocol\_amfr.md}，哈希 3246ffda，2026-11-05T01:30:37Z 锁定，之后才求解任何测试算例。

\paragraph{测试算例（锁定前从未求解）。}
T-A：gen3t，$m=64$，$F\in\{10^{-8},10^{-4},3\times10^{-2},0.3,3\}\times G\in\{10^{-3},3\times10^{-2},0.3,10\}$，另加 $m=128$ 两个（共 22 个，最大 630 万未知量）。
T-B：新写的多群辐射扩散 + 物质温度（\texttt{mgrd.py}，对称 M 矩阵）：$n_g\in\{2,4,8\}$ 群，不透明度 $\bar\sigma\in\{10^{-2},1,10^2\}$，物质热容 $\zeta\in\{10^{-2},10^2\}$，物质热传导 $\lambda\in\{0,1\}$，种子 7，$m=40$（共 36 个，最多 9 个场、58 万未知量）。
T-C：真实 SC26\_01。方法：amg-3d（整体 BoomerAMG）、amg-unk（hypre systems AMG，unknown 方法）、amfr，以及消融 amfr-sep（不合并）、amfr-noelim（不消去）、amfr-all（消去后全部合并）。
均为 GMRES(50)、右预条件、求解器自报残差 $10^{-8}$、$x_0=0$、最多 500 步、单线程、顺序运行。

\section{结果}
\begin{figure}[H]
\centering
\includegraphics[width=\linewidth]{ratios.png}
\caption{各方法总时间相对 amg-3d 的比值（对数轴，低于 1 表示更快）。蓝色圆点为 AMFR。}
\end{figure}

\begin{table}[H]
\centering\footnotesize
\begin{tabular}{rll rr l rr l}
\toprule
$m$ & $F$ & $G$ & amg-3d & amg-unk & AMFR 结构 & AMFR & 比值 & 最快消融 \\
\midrule
64 & 1e-8 & 1e-3 & 18 / 1.77 & 31 / 2.43 & \texttt{e1|2|0} & 18 / 1.51 & 0.85 & amfr \\
64 & 1e-8 & 3e-2 & 18 / 1.70 & 47 / 3.34 & \texttt{|12|0} & 18 / 1.97 & 1.16 & noelim \\
64 & 1e-8 & 0.3 & 18 / 2.14 & 47 / 4.31 & \texttt{|12|0} & 18 / 2.13 & 1.00 & amfr \\
64 & 1e-8 & 10 & 19 / 3.08 & 33 / 3.78 & \texttt{|12|0} & 18 / 2.42 & 0.79 & noelim \\
64 & 1e-4 & 1e-3 & 18 / 1.64 & 31 / 2.33 & \texttt{e1|2|0} & 18 / 1.49 & 0.91 & sep \\
64 & 1e-4 & 3e-2 & 18 / 1.66 & 46 / 3.24 & \texttt{|12|0} & 18 / 1.80 & 1.08 & amfr \\
64 & 1e-4 & 0.3 & 18 / 2.17 & 47 / 4.21 & \texttt{|12|0} & 18 / 2.05 & 0.94 & amfr \\
64 & 1e-4 & 10 & 18 / 2.71 & 33 / 3.93 & \texttt{|12|0} & 18 / 2.50 & 0.92 & amfr \\
64 & 3e-2 & 1e-3 & 17 / 1.61 & 30 / 2.34 & \texttt{e1|2|0} & 17 / 1.39 & 0.86 & amfr \\
64 & 3e-2 & 3e-2 & 18 / 1.81 & 34 / 2.56 & \texttt{|12|0} & 17 / 2.14 & 1.19 & noelim \\
64 & 3e-2 & 0.3 & 17 / 2.34 & 33 / 3.16 & \texttt{|12|0} & 17 / 2.05 & 0.87 & noelim \\
64 & 3e-2 & 10 & 18 / 2.55 & 29 / 3.47 & \texttt{|12|0} & 18 / 2.71 & 1.07 & noelim \\
64 & 0.3 & 1e-3 & 16 / 1.75 & 18 / 1.58 & \texttt{e1|02} & 16 / 1.68 & 0.96 & sep \\
64 & 0.3 & 3e-2 & 16 / 1.80 & 18 / 1.63 & \texttt{|12|0} & 16 / 1.69 & 0.94 & noelim \\
64 & 0.3 & 0.3 & 16 / 1.79 & 18 / 2.15 & \texttt{|12|0} & 17 / 1.99 & 1.11 & noelim \\
64 & 0.3 & 10 & 16 / 2.28 & 19 / 2.71 & \texttt{|12|0} & 17 / 2.26 & 0.99 & sep \\
64 & 3 & 1e-3 & 15 / 1.90 & 15 / 1.44 & \texttt{e1|02} & 15 / 1.63 & 0.86 & sep \\
64 & 3 & 3e-2 & 15 / 1.86 & 15 / 1.46 & \texttt{|01|2} & 15 / 1.65 & 0.89 & sep \\
64 & 3 & 0.3 & 15 / 1.83 & 15 / 1.94 & \texttt{|01|2} & 15 / 1.61 & 0.88 & noelim \\
64 & 3 & 10 & 17 / 2.33 & 16 / 2.35 & \texttt{|01|2} & 17 / 2.17 & 0.93 & sep \\
128 & 1e-8 & 1e-3 & 21 / 21.26 & 32 / 29.50 & \texttt{e1|2|0} & 22 / 16.59 & 0.78 & sep \\
128 & 1e-8 & 0.3 & 21 / 22.58 & 49 / 56.10 & \texttt{|12|0} & 20 / 22.30 & 0.99 & noelim \\
\bottomrule
\end{tabular}
\caption{T-A（三温）。单元格为“迭代 / 总时间 (s)”；结构记号：\texttt{e1} 表示消去场 1（电子），\texttt{|12|0} 表示先联合解 \{1,2\}、再解 \{0\}。场 0 = 辐射，1 = 电子，2 = 离子。}
\end{table}

{\footnotesize\setlength{\tabcolsep}{4pt}
\begin{longtable}{rlll rr l rr l}
\caption{T-B（多群 + 物质温度；场 $0\ldots n_g-1$ 为各群，$n_g$ 为物质温度）。}\\
\toprule
$n_g$ & $\bar\sigma$ & $\zeta$ & $\lambda$ & amg-3d & amg-unk & AMFR 结构 & AMFR & 比值 & 最快消融 \\
\midrule
\endhead
2 & 1e-2 & 1e-2 & 0 & 83 / 1.12 & 83 / 0.98 & \texttt{e2|0|1} & 14 / 0.29 & 0.26 & sep \\
2 & 1e-2 & 1e-2 & 1 & 33 / 0.71 & 44 / 0.84 & \texttt{|1|0|2} & 83 / 1.38 & 1.93 & all \\
2 & 1e-2 & 1e2 & 0 & 11 / 0.27 & 11 / 0.23 & \texttt{e2|0|1} & 11 / 0.24 & 0.91 & noelim \\
2 & 1e-2 & 1e2 & 1 & 11 / 0.41 & 11 / 0.35 & \texttt{|1|2|0} & 11 / 0.32 & 0.80 & amfr \\
2 & 1 & 1e-2 & 0 & 467 / 4.98 & 468 / 4.43 & \texttt{e2|0|1} & 40 / 0.56 & 0.11 & all \\
2 & 1 & 1e-2 & 1 & 46 / 1.45 & 101 / 1.44 & \texttt{|0|2|1} & 225 / 3.25 & 2.25 & all \\
2 & 1 & 1e2 & 0 & 8 / 0.16 & 9 / 0.18 & \texttt{e2|0|1} & 8 / 0.18 & 1.13 & sep \\
2 & 1 & 1e2 & 1 & 8 / 0.29 & 9 / 0.27 & \texttt{|2|0|1} & 10 / 0.29 & 1.01 & sep \\
2 & 1e2 & 1e-2 & 0 & 失败 & 失败 & \texttt{e2e0|1} & 8 / 0.12 & 0.02 & amfr \\
2 & 1e2 & 1e-2 & 1 & 45 / 1.37 & 198 / 2.42 & \texttt{e0|12} & 11 / 0.29 & 0.21 & all \\
2 & 1e2 & 1e2 & 0 & 4 / 0.11 & 7 / 0.14 & \texttt{e2e0|1} & 4 / 0.10 & 0.95 & amfr \\
2 & 1e2 & 1e2 & 1 & 5 / 0.23 & 6 / 0.21 & \texttt{e0|2|1} & 5 / 0.19 & 0.86 & sep \\
4 & 1e-2 & 1e-2 & 0 & 16 / 0.58 & 16 / 0.57 & \texttt{e4|3|0|1|2} & 15 / 0.62 & 1.06 & sep \\
4 & 1e-2 & 1e-2 & 1 & 16 / 0.61 & 15 / 0.77 & \texttt{|3|0|1|2|4} & 20 / 0.77 & 1.27 & all \\
4 & 1e-2 & 1e2 & 0 & 11 / 0.46 & 11 / 0.47 & \texttt{e4|3|0|1|2} & 11 / 0.50 & 1.07 & noelim \\
4 & 1e-2 & 1e2 & 1 & 11 / 0.66 & 11 / 0.67 & \texttt{|3|4|0|1|2} & 11 / 0.56 & 0.85 & amfr \\
4 & 1 & 1e-2 & 0 & 199 / 4.60 & 199 / 3.86 & \texttt{e4|0|3|1|2} & 182 / 4.22 & 0.92 & all \\
4 & 1 & 1e-2 & 1 & 63 / 1.71 & 52 / 1.55 & \texttt{|0|3|1|4|2} & 108 / 2.85 & 1.67 & all \\
4 & 1 & 1e2 & 0 & 8 / 0.35 & 8 / 0.36 & \texttt{e4|0|1|2|3} & 8 / 0.42 & 1.20 & noelim \\
4 & 1 & 1e2 & 1 & 8 / 0.52 & 8 / 0.51 & \texttt{|4|0|1|2|3} & 8 / 0.47 & 0.91 & noelim \\
4 & 1e2 & 1e-2 & 0 & 241 / 4.78 & 241 / 3.73 & \texttt{e4e0e1|3|2} & 13 / 0.36 & 0.07 & sep \\
4 & 1e2 & 1e-2 & 1 & 90 / 2.07 & 79 / 1.86 & \texttt{e0e1|24|3} & 14 / 0.63 & 0.30 & amfr \\
4 & 1e2 & 1e2 & 0 & 3 / 0.16 & 4 / 0.20 & \texttt{e4e0e1|2|3} & 3 / 0.22 & 1.36 & sep \\
4 & 1e2 & 1e2 & 1 & 4 / 0.30 & 4 / 0.28 & \texttt{e0e1|4|2|3} & 5 / 0.27 & 0.89 & amfr \\
8 & 1e-2 & 1e-2 & 0 & 17 / 1.24 & 17 / 1.33 & \texttt{e8|7|0|6|1|2|3|4|5} & 17 / 1.48 & 1.20 & noelim \\
8 & 1e-2 & 1e-2 & 1 & 17 / 1.27 & 15 / 1.55 & \texttt{|7|0|6|1|2|5|3|4|8} & 21 / 1.60 & 1.25 & sep \\
8 & 1e-2 & 1e2 & 0 & 11 / 1.08 & 11 / 1.07 & \texttt{e8|7|0|1|6|2|3|4|5} & 11 / 1.21 & 1.12 & noelim \\
8 & 1e-2 & 1e2 & 1 & 11 / 1.35 & 11 / 1.38 & \texttt{|7|0|8|1|2|3|4|5|6} & 11 / 1.09 & 0.81 & sep \\
8 & 1 & 1e-2 & 0 & 194 / 9.34 & 194 / 8.09 & \texttt{e8|0|1|7|2|3|4|5|6} & 205 / 9.26 & 0.99 & amfr \\
8 & 1 & 1e-2 & 1 & 62 / 3.34 & 51 / 3.24 & \texttt{|0|7|1|2|3|6|4|8|5} & 104 / 5.78 & 1.73 & all \\
8 & 1 & 1e2 & 0 & 8 / 0.76 & 8 / 0.84 & \texttt{e8|0|1|7|2|3|4|5|6} & 8 / 1.10 & 1.44 & noelim \\
8 & 1 & 1e2 & 1 & 9 / 1.11 & 8 / 1.12 & \texttt{|0|8|1|2|3|4|5|6|7} & 8 / 1.04 & 0.94 & noelim \\
8 & 1e2 & 1e-2 & 0 & 239 / 9.75 & 239 / 8.10 & \texttt{e8e0e1e2|7|3|4|5|6} & 159 / 5.32 & 0.55 & all \\
8 & 1e2 & 1e-2 & 1 & 90 / 4.03 & 79 / 3.05 & \texttt{e0e1e2|7|3|6|4|8|5} & 96 / 3.55 & 0.88 & all \\
8 & 1e2 & 1e2 & 0 & 3 / 0.30 & 4 / 0.40 & \texttt{e8e0e1e2|7|3|4|5|6} & 3 / 0.64 & 2.12 & noelim \\
8 & 1e2 & 1e2 & 1 & 4 / 0.52 & 4 / 0.51 & \texttt{e0e1e2|8|3|4|5|6|7} & 5 / 0.54 & 1.03 & amfr \\
\bottomrule
\end{longtable}}

\paragraph{真实 SC26\_01（同批次）。} AMFR 自动消去场 1（电子，非对角占比 $2.1\times10^{-4}$），离子 → 辐射，即 bt-v：23 步，16.2\,s（设置 6.2\,s）；amg-3d 23 步，18.1\,s（设置 5.9\,s）。

\paragraph{要点。}
\begin{itemize}
\item \textbf{三温：} AMFR 与 bt-v / bt-j 一致；与 BoomerAMG 迭代数相同，时间中位数 0.93（$m=128$：0.78、0.99）。消融 amfr-sep 在 $G\ge3\times10^{-2}$ 时 65--72 步（$m=128$ 慢 5.2 倍），说明合并判据在这里是必要的、并且判对了。
\item \textbf{多群，BoomerAMG 弱的算例}（amg-3d $\ge40$ 步的 13 个）：AMFR 的时间比为 0.02、0.07、0.11、0.21、0.26、0.30、0.55，收益来自消去物质温度或光学厚的群（例如 500 步不收敛 → 8 步）；不消去的消融（amfr-noelim）在这些算例上与 BoomerAMG 一样差或更差。
\item \textbf{失败类：} $\lambda=1$（物质温度有热传导，不能消去）且 $\zeta=10^{-2}$（热容小，物质温度与各群强耦合）：AMFR 把物质温度单独解，83/225/108/104 步，对 amg-3d 33/46/63/62 步（1.67--2.25 倍时间）；amfr-all 与 BoomerAMG 持平。
\item \textbf{开销：} 迭代很少（3--17 步）的多场算例上，AMFR 慢 1.1--2.1 倍（绝对差 0.1--0.6\,s）：Python 中逐场的 V 循环和耦合乘法开销，以及消去星形中心后产生的稠密单元耦合（amfr-all 在 8 群时 2.4--2.7 倍）。
\item \textbf{hypre systems AMG（amg-unk）}在三温问题上比 amg-3d 慢（中位数 1.36），在多群上持平（0.98）；AMFR 只在 33/58 个算例上快于它。
\end{itemize}

\section{预测检验}
\begin{enumerate}
\item 58 个全部收敛，且至少 55 个迭代不超过 amg-3d 的 1.5 倍：\textbf{不成立}（58/58 收敛，54/58；四个例外都是上述失败类）。
\item 时间比中位数 T-A $\le0.90$、T-B $\le0.75$，且至少 52/58 不超过 1.10：\textbf{不成立}（0.93、0.97；42/58）。
\item 至少 49/58 的 AMFR 时间不超过最快消融的 1.15 倍：\textbf{不成立}（45/58）。
\item 至少 46/58 快于 amg-unk：\textbf{不成立}（33/58）。
\item 真实矩阵上 AMFR 选出 \texttt{e1|2|0} 并快于 amg-3d：\textbf{成立}（16.2\,s 对 18.1\,s）。
\item T-A $m=64$：$G=10^{-3}$ 的 5 个全部消去电子，$G\ge3\times10^{-2}$ 的 15 个全不消去：\textbf{成立}（5/5，0/15）。
\end{enumerate}

\section{锁定后的诊断（探索性，未锁定；开发种子 1，$m=20$--32）}
\begin{itemize}
\item \textbf{失败原因是判据，不是 V 循环。} 在失败类算例（2 群，$\bar\sigma=1$，$\zeta=10^{-2}$，$\lambda=1$）上把每组的 V 循环换成精确 LU：分开解仍要 159 步，合并解 1 步。所以耦合确实很强，是判据把它判弱了。
\item \textbf{截断是 bug。} 物质温度行不是对角占优：$\sum_g\sigma_g\sqrt{w_g}>c_v+\sum_g\sigma_g w_g$，耦合/行和 $\approx2.5$。H50 对单边耦合 $b_{ij}/\rho_i$ 截断到 1，回路增益被压到 $0.39<0.7$，实际约为 1。
而且只有乘积 $c_{ij}c_{ji}$（以及整个耦合矩阵的谱半径）在场的单位变换下不变（单边耦合按 $t_j/t_i$ 变），所以截断单边耦合本身就不符合方法的设计原则。
\item \textbf{修正（H51 候选）：} 不截断；每个单元组装常数模态耦合矩阵 $C_{ij}=b_{ij}/\rho_i$，去掉组内元素后若谱半径平方 $>\tau_g$ 的单元比例 $\ge\varphi$，就合并贡献最大的两组，直到不再满足。两个场时与 H50 的判据完全相同；也能识别“许多弱耦合加起来很强”的星形耦合。
在 18 个开发算例上：失败类的迭代降到与 BoomerAMG 相同（2 群：68 → 32、155 → 49；4 群：88 → 51；8 群：87 → 50），但时间只是持平（0.90--1.42），因为合并后基本就是整体 AMG 加 Python 开销；8 群光学厚的算例中 amfr-all 仍明显更好（8 步对 51 步），说明多场时判据仍偏向分开。
\end{itemize}

\section{结论}
\begin{enumerate}
\item 自动多场约化是可行的：同一个只看矩阵的规则在三温问题上自动给出 bt-v / bt-j，在真实 SC26\_01 上给出 bt-v；三温问题上最多比 BoomerAMG 慢 1.19 倍。在多群问题上则有一类算例慢到 2.25 倍（见第 3 点）。
\item 真正的大收益来自\textbf{消去}：当某个场（无扩散的物质温度、光学厚的群）的块接近对角时，消去后剩下的问题对 AMG 容易得多，例如 500 步不收敛 → 8 步。这一点不依赖三温结构，在多群问题上同样成立。
\item \textbf{合并}判据决定的是“分开解省一点（0.8--0.9 倍）”还是“分开错了亏很多（2 倍以上）”，风险不对称。H50 的判据有截断 bug；修正后不再亏迭代，但也不再省时间。
\item 局限：Python 实现中逐场开销在简单算例上明显；消去星形中心会产生稠密单元耦合；所有比较都是单线程。
\end{enumerate}

\paragraph{下一步（需要决定）。}
(a) 把修正后的合并判据在新种子、新参数上锁定验证（H51）；(b) 在消去规则里加“填充”判据（避免消去星形中心后单元耦合变稠密）；(c) 把 AMFR 写进 C 版本（与 bt-v 的 C 实现共用），去掉 Python 开销后再比时间。

\paragraph{文件。} \sloppy 代码 \texttt{FD\_LGF/multifield/}：\texttt{amfr.py}（方法；\texttt{join="spectral"} 为 H51 探索版本）、\texttt{mgrd.py}（多群问题）、\texttt{run\_mf.py}、\texttt{test.sh}、\texttt{analyze.py}；结果 \texttt{results/test\_*.csv}，开发数据 \texttt{results/dev*\_*.csv}。
\end{document}
